\documentclass[10pt,a4paper]{amsart}

\pdfinfoomitdate=1
\pdftrailerid{}
\pdfsuppressptexinfo=15

\usepackage[T1]{fontenc}
\usepackage{lmodern}
\usepackage[margin=25mm]{geometry}
\usepackage[protrusion=true,expansion=false]{microtype}
\usepackage{amsmath,amssymb,mathtools}
\usepackage{booktabs}
\usepackage{enumitem}
\usepackage{xcolor}
\usepackage[numbers,sort&compress]{natbib}
\usepackage[colorlinks=true,linkcolor=blue!55!black,citecolor=blue!55!black,urlcolor=blue!55!black]{hyperref}
\usepackage[nameinlink,noabbrev]{cleveref}

\hypersetup{pdfauthor={},pdftitle={},pdfsubject={},pdfkeywords={}}
\setlist{nosep,leftmargin=*}
\raggedbottom
\allowdisplaybreaks
\emergencystretch=2em

\newtheorem{theorem}{Theorem}[section]
\newtheorem{proposition}[theorem]{Proposition}
\newtheorem{lemma}[theorem]{Lemma}
\newtheorem{corollary}[theorem]{Corollary}
\theoremstyle{definition}
\newtheorem{definition}[theorem]{Definition}
\theoremstyle{remark}
\newtheorem{remark}[theorem]{Remark}

\newcommand{\Bn}{\{0,1\}^n}
\newcommand{\Lyn}{\mathcal L}
\newcommand{\depth}{\operatorname{depth}}
\newcommand{\im}{\operatorname{im}}
\newcommand{\ones}{\mathbf 1}
\newcommand{\ind}{\mathbf 1}

\title[Lyndon-factor-start feedback]{Iterating Lyndon-Factor Starts on Binary Words}
\author{Anonymous}
\date{}
\subjclass[2020]{37P25, 68R15, 05A15}
\keywords{finite dynamical system, Lyndon word, Chen--Fox--Lyndon
factorization, suffix minimum, fibre}

\begin{document}

\begin{abstract}
Fix $0<1$ and replace a binary word by the mask of the starting positions in
its nonincreasing Chen--Fox--Lyndon factorization, and iterate this mask map.
The classical suffix-array characterization of factor starts is imported as
owned static input and receives no contribution credit.  In the resulting
dynamics, every leading $1$ is a singleton Lyndon factor, so the leading-one
prefix grows under every nonfixed update.  We prove that $1^n$ is the unique
recurrent state and that every orbit reaches it within $n$ steps.  The bound
is sharp at a unique state, the alternating word $0101\cdots$.  For every
target mask, we also identify the complete fibre with the nonincreasing
chains of binary Lyndon words having prescribed lengths and express its size
as a product of rectangular comparison matrices.  The one-factor and
all-singleton targets give the classical binary Lyndon census and $n+1$.
All static factorization, suffix-record, and enumeration results are credited
background.  External release remains on hold.
\end{abstract}

\maketitle

\section{Factor-start feedback and the credit boundary}

Words are ordered lexicographically over $0<1$, with a proper prefix smaller
than the longer word.  A nonempty word $u$ is \emph{Lyndon} if it is strictly
smaller than each of its nontrivial rotations.  The classical
Chen--Fox--Lyndon theorem \cite{ChenFoxLyndon1958} gives every nonempty word a
unique factorization
\begin{equation}\label{eq:CFL}
 w=u_1u_2\cdots u_k,\qquad u_1\ge u_2\ge\cdots\ge u_k,
\end{equation}
into Lyndon words; repeated factors are written separately.

Let $1=s_1<s_2<\cdots<s_k\le n$ be the factor starts in \eqref{eq:CFL}.  Define
the binary self-map $L_n:\Bn\to\Bn$ by
\begin{equation}\label{eq:map}
 L_n(w)_i=1\quad\Longleftrightarrow\quad i\in\{s_1,\ldots,s_k\}.
\end{equation}
The factorization is recomputed from the mask after every update.

Duval computes \eqref{eq:CFL}, the least suffix, and the least circular shift
in linear time \cite{Duval1983}.  Lyndon arrays and forests have their own
algorithmic theory \cite{FranekIslamRahmanSmyth2016,BadkobehCrochemore2022}.
Mantaci, Restivo, Rosone, and Sciortino prove that the CFL factor starts are
exactly the left-to-right minima of the suffix-rank permutation
\cite[Theorem~2.2]{MantaciRestivoRosoneSciortino2014}.  The equivalent strict
suffix-record formulation and the ordered-tail comparison used to recover it
are therefore owned static inputs.  We restate them in \cref{sec:owned} only
to fix conventions; they receive zero contribution credit.

The CFL theorem, all suffix-array and factor-start characterizations, the
algorithms above, the binary Lyndon census, necklace M\"obius inversion, and
matrix multiplication are likewise zero-credit background.  After this
subtraction, the residual package consists only of the recurrent dynamics of
the repeated map \eqref{eq:map}, its sharp unique depth-$n$ state, and its
target-wise ordered-Lyndon fibre atlas.

A state is recurrent if it lies on a directed cycle.  Its depth is the first
time its orbit reaches a recurrent state.

\section{Owned static interface: suffix records}\label{sec:owned}

The suffix-record equivalence in this section is exactly the static result of
\citet[Theorem~2.2]{MantaciRestivoRosoneSciortino2014}, written in suffix
notation rather than inverse-suffix-array notation.  We reproduce a proof so
that the orientation and strictness convention are auditable, but neither the
statement nor the proof is a residual contribution.

The proof uses two classical consequences of the Lyndon definition.  A Lyndon
word is unbordered and is strictly smaller than each of its nonempty proper
suffixes.  It also uses the ordered-tail comparison underlying the CFL
theorem: if
$u_1\ge\cdots\ge u_k$ are Lyndon and
$S_h=u_h\cdots u_k$, then
\begin{equation}\label{eq:taildescent}
 S_1>S_2>\cdots>S_k.
\end{equation}
Here is the short reduction, including equal factors.  It is enough to
compare adjacent tails.  Fix $h$, put $u=u_h$, and remove from both tails the
common initial copies of $u$ occurring in $S_{h+1}$.  If no word remains on
the right, that right-hand word is a proper prefix.  Otherwise it begins with
a Lyndon factor $v<u$; the classical Lyndon concatenation comparison says
$uz>z$ whenever the CFL factors of $z$ begin with $v<u$ and thereafter do
not increase.  Thus $S_h>S_{h+1}$ in both cases.  This is a static comparison
lemma, not a dynamical contribution; it is also the comparison step in the
standard proof of the CFL factorization \cite{ChenFoxLyndon1958,Duval1983}.
This ordered-tail comparison is owned static machinery and receives zero
contribution credit as well.

\begin{proposition}[Mantaci--Restivo--Rosone--Sciortino; owned input]
\label{prop:records}
For $w=w_1\cdots w_n$, position $i$ starts a factor in \eqref{eq:CFL} if and
only if its suffix is a strict new minimum among the suffixes seen from the
left:
\begin{equation}\label{eq:record}
 L_n(w)_i=1
 \quad\Longleftrightarrow\quad
 w_i\cdots w_n<w_j\cdots w_n\quad\text{for every }j<i.
\end{equation}
\end{proposition}

\begin{proof}[Reproduced background proof]
Write $S_h=u_hu_{h+1}\cdots u_k$ for the suffix beginning at factor $u_h$.
By \eqref{eq:taildescent}, each later factor-start suffix is strictly smaller
than every earlier factor-start suffix; repeated equal factors are included
in the proper-prefix case of the reduction above.

Now take a position $q$ strictly inside $u_h$, and write $v$ for the proper
suffix of $u_h$ beginning there.  Since $u_h$ is Lyndon, $u_h<v$.  The two
words are not in a prefix relation: $v$ is shorter, while $v$ cannot be a
prefix of $u_h$ because that would make it a border.  Appending the same later
factor tail therefore preserves the strict comparison, so $S_h$ is smaller
than the suffix beginning at $q$.

If $i=s_j$ is a factor start, every earlier position is either an earlier
factor start, whose suffix is larger by the first paragraph, or lies inside
an earlier factor, whose suffix is larger than that factor-start suffix by the
second paragraph.  Hence \eqref{eq:record} holds.  Conversely, if $i$ lies
inside $u_h$, the earlier position $s_h$ has a strictly smaller suffix, so
$i$ is not a new minimum.  This proves both directions.
\end{proof}

\begin{remark}
Strictness matters when consecutive factors are equal.  For example, the CFL
factorization of $00$ is $0\,0$, and the second suffix $0$ is strictly smaller
than the first suffix $00$ because it is a proper prefix.
\end{remark}

Thus \cref{prop:records} supplies an intrinsic implementation check for the
map, but it is not counted among the paper's residual theorems.

\section{The leading-one amplifier}

\begin{lemma}[literal prefix rule]\label{lem:prefix}
If $w=1^r0s$ with $0\le r<n$, then
\begin{equation}\label{eq:prefix}
 L_n(w)=1^r L_{n-r}(0s).
\end{equation}
More generally, $L_{m+1}(1v)=1L_m(v)$ for every $v\in\{0,1\}^m$.
\end{lemma}

\begin{proof}
A binary Lyndon word of length greater than one begins with zero: a word
beginning in one and containing a zero has a smaller rotation beginning at
that zero, while $1^d$ is not primitive for $d>1$.  Hence every leading one
in $w$ is a singleton Lyndon factor.  The first factor of the nonempty
remainder $0s$ begins in zero and is therefore smaller than the singleton
$1$.  Uniqueness of \eqref{eq:CFL} shows that the full factorization is the
$r$ singleton factors followed by the CFL factorization of $0s$, proving
\eqref{eq:prefix}.

For arbitrary $v$, either $v=1^m$, when the last statement is immediate, or
$v=1^q0s$ and \eqref{eq:prefix} applied with $r=q$ and $r=q+1$ gives the
identity.
\end{proof}

Let $\lambda(w)$ be the number of leading ones in $w$.

\begin{theorem}[recurrent atlas and clock bound]\label{thm:atlas}
The word $1^n$ is the unique fixed point and the unique recurrent state of
$L_n$.  Every orbit reaches it in at most $n$ updates.
\end{theorem}

\begin{proof}
The CFL factorization of $1^n$ consists of $n$ singleton factors, so $1^n$ is
fixed.  If $w\ne1^n$, write $w=1^r0s$.  Equation \eqref{eq:prefix} shows that
the first $r$ output bits are one, and the first bit of
$L_{n-r}(0s)$ is also one because every nonempty factorization has a first
factor.  Thus $\lambda(L_n(w))\ge r+1$.  The integer $\lambda$ strictly
increases at every nonfixed update and is bounded by $n$.  No other recurrent
state can exist, and at most $n$ updates are required.
\end{proof}

\section{The sharp and unique deepest word}

Let $a_n=0101\cdots$ be the length-$n$ alternating word beginning in zero.

\begin{lemma}[alternating recursion]\label{lem:alternating}
For every $n\ge1$,
\begin{equation}\label{eq:alternating}
 L_n(a_n)=1a_{n-1},
\end{equation}
where $a_0$ is empty.
\end{lemma}

\begin{proof}
The word $01$ is Lyndon.  If $n=2m$, the CFL factorization of $a_n$ is
$(01)^m$, written as $m$ equal factors.  If $n=2m+1$, it is
$(01)^m0$; this is nonincreasing because $01>0$.  The starts are exactly the
odd positions, which is the mask $1a_{n-1}$.  At $n=1$, the factorization is
the singleton $0$ and the same identity holds.
\end{proof}

\begin{theorem}[sharp clock and uniqueness]\label{thm:unique}
The maximum depth of $L_n$ is $n$.  Its unique depth-$n$ state is $a_n$.
\end{theorem}

\begin{proof}
By \cref{lem:prefix}, prefixing a word by one does not change its depth:
the identity $L_{m+1}(1v)=1L_m(v)$ persists under iteration, including when
$v$ becomes all ones.  Therefore \eqref{eq:alternating} and induction give
\[
 \depth(a_n)=1+\depth(a_{n-1})=n,
\]
with $a_1=0\mapsto1$.  This attains the upper bound in \cref{thm:atlas}.

We prove uniqueness by induction.  At $n=1$, the word $0$ is the only state
of depth one.  Let $n\ge2$ and suppose $w$ has depth $n$.  Its image has depth
$n-1$ and begins in one.  Write $L_n(w)=1v$.  Prefix invariance of depth and
the induction hypothesis force $v=a_{n-1}$, so
\begin{equation}\label{eq:forcedmask}
 L_n(w)=1a_{n-1}.
\end{equation}
The one-positions in this mask prescribe CFL factor lengths
$(2,2,\ldots,2)$ when $n$ is even and $(2,2,\ldots,2,1)$ when $n$ is odd.
The only binary Lyndon word of length two is $01$.  In the odd case, the last
singleton cannot be $1$, because the required nonincrease would read
$01\ge1$, which is false; it must be $0$.  Hence $w=(01)^{n/2}$ for even
$n$, and $w=(01)^{(n-1)/2}0$ for odd $n$.  In both cases $w=a_n$.
\end{proof}

\section{Every target fibre as an ordered Lyndon chain}

For $d\ge1$, let $\Lyn_d$ be the lexicographically ordered set of binary
Lyndon words of length $d$.  If a target $y\in\Bn$ begins in one, its
one-positions determine a composition
$\ell(y)=(\ell_1,\ldots,\ell_k)$ of $n$: these are the successive gaps
between factor starts, with the last gap ending at $n+1$.

\begin{theorem}[complete inverse atlas]\label{thm:fibre}
If $y_1=0$, then $L_n^{-1}(y)$ is empty.  If $y_1=1$ and
$\ell(y)=(\ell_1,\ldots,\ell_k)$, then
\begin{equation}\label{eq:fibre}
 |L_n^{-1}(y)|
 =\#\{(u_1,\ldots,u_k):u_i\in\Lyn_{\ell_i},\ 
                         u_1\ge\cdots\ge u_k\}.
\end{equation}
Thus positivity of the right-hand side is the complete image criterion.
\end{theorem}

\par\smallskip
\begin{proof}
Every factor-start mask begins with one, giving the first assertion.  For a
target beginning in one, any preimage has factors starting at exactly its
one-positions, so their lengths are $\ell_1,\ldots,\ell_k$.  By the CFL
theorem, those factors are Lyndon and nonincreasing, which maps the fibre into
the set on the right of \eqref{eq:fibre}.  Conversely, concatenate any tuple
on the right.  It is already a nonincreasing factorization into Lyndon words;
uniqueness of \eqref{eq:CFL} makes it the CFL factorization of the
concatenation, whose start mask is $y$.  The two constructions are inverse.
\end{proof}

\par\smallskip
The count has a finite matrix form.  For positive $a,b$, let
$M_{a,b}$ be the $|\Lyn_a|\times|\Lyn_b|$ zero-one matrix
\begin{equation}\label{eq:matrix}
 M_{a,b}(u,v)=\ind\{u\ge v\},\qquad u\in\Lyn_a, v\in\Lyn_b,
\end{equation}
and let $\mathbf 1_d$ be the all-one column indexed by $\Lyn_d$.

\par\smallskip
\begin{corollary}[matrix fibres and two closed cells]\label{cor:special}
For $k\ge2$, the fibre in \eqref{eq:fibre} equals
\begin{equation}\label{eq:product}
 \mathbf 1_{\ell_1}^{\mathsf T}
 M_{\ell_1,\ell_2}\cdots
 M_{\ell_{k-1},\ell_k}\mathbf 1_{\ell_k};
\end{equation}
for $k=1$ it equals $|\Lyn_{\ell_1}|$.  Moreover,
\begin{align}
 |L_n^{-1}(10^{n-1})|
   &=\frac1n\sum_{d\mid n}\mu(d)2^{n/d},\label{eq:onefactor}\\
 |L_n^{-1}(1^n)|&=n+1.\label{eq:allstarts}
\end{align}
\end{corollary}

\begin{proof}
Expanding \eqref{eq:product} sums one for each tuple satisfying all adjacent
inequalities, exactly as in \eqref{eq:fibre}.  The target $10^{n-1}$
prescribes the one-part composition $(n)$, so its fibre is $\Lyn_n$.  The
classical primitive-necklace M\"obius formula gives \eqref{eq:onefactor}.
For $1^n$, every factor has length one.  The two one-letter Lyndon words are
$0$ and $1$, and a nonincreasing sequence is uniquely $1^j0^{n-j}$ for some
$0\le j\le n$, proving \eqref{eq:allstarts}.
\end{proof}

\section{Exact control and limitations}

The paper-local verifier independently computes the Duval mask and the strict
suffix-record mask for every binary word through $n=18$.  Their comparison is
an integration check for the imported theorem of Mantaci et al., not evidence
for a residual claim.  The verifier also constructs every functional graph in
that range, totaling $524{,}286$ states, and checks the fixed atlas, clock, and
unique deepest word.  A separate Lyndon-word generator and rectangular-matrix
evaluator agrees with every literal target fibre through $n=14$; the two
special fibres are checked through $n=18$.  The canonical run executes
$2{,}654{,}300$ exact assertions.

Finite enumeration is only counterexample pressure.  The observed image and
largest-fibre sequences are not promoted to unproved all-length statements.
The suffix-record theorem and ordered-tail comparison are explicitly
owner-subtracted; the residual is limited to the iterated start-mask dynamics,
the sharp unique depth-$n$ orbit, and the ordered-Lyndon fibre atlas.  The
owner search for that residual is bounded, so a non-hit is not a novelty or
priority certificate.  This anonymous artifact remains
\texttt{HOLD\_EXTERNAL}; public posting, submission, and release are not
authorized.

\bibliographystyle{plainnat}
\bibliography{references}

\end{document}
